\BourbakiExercisesSection

\begin{enumerate}
\def\labelenumi{\arabic{enumi}.}
\item
  若 V 是有限生成投射 A-模，则 End(V) 亦然，且典范地等同于 \(V \otimes_{A} V^{*}\) （II，§4，第 2 号，命题 2 的推论）。若将每个元素 \(u \in \text{End}(V)\) 与 End(V) 上的 A-线性形式 \(v \mapsto \text{Tr}(uv)\) 相关联，就定义了从 End(V) 到其对偶 (End(V))* 上的一个 A-线性双射。在此映射下将 End(V) 与其对偶 A-模等同，则 End(V) 上的 A-代数结构通过对偶性（第 1 号，例 3）在 End(V) 上定义一个 A-余代数结构，它是余结合的，并以映为 \(V \to A\) 的形式 Tr 作为余单位。特别地，对 \(V = A^{n}\) ，在 \(\mathbf{M}_{n}(\mathbf{A})\) 上定义了一个余代数结构，其余积由 \(c(E_{ij}) = \sum_{k} E_{kj} \otimes E_{ik}\) 给出。这个余代数结构与 \(\mathbf{M}_{n}(\mathbf{A})\) 上的代数结构不定义双代数结构。
\item
  证明在定义 3（第 4 号）中，条件（3）（关于分次双代数的）等价于：乘积 \(m: \mathbf{E} \otimes \mathbf{E} \to \mathbf{E}\) 是从分次余代数 \(\mathbf{E} \otimes \mathbf{E}\) 到分次余代数 \(\mathbf{E}\) 中的态射。
\item
  证明：若 \(\mathbf{M}\) 是一个 A-模，则在 \(\mathsf{T}(\mathbf{M})\) 上存在且仅存在一个余交换双代数结构，其代数结构是通常的结构，且其余积 \(c\) 满足：对所有 \(x\in \mathbf{M}\) ， \(c(x) = 1\otimes x + x\otimes 1\) 。
\item
  设 \(\mathbf{E}\) 是 \(\mathbf{A}\) 上的一个交换双代数， \(m\) 是乘积 \(\mathbf{E} \otimes \mathbf{E} \to \mathbf{E}\) ， \(c\) 是余积 \(\mathbf{E} \to \mathbf{E} \otimes \mathbf{E}\) ， \(e\) 是单位元，而 \(\gamma\) 是 \(\mathbf{E}\) 的余单位。
\end{enumerate}

\begin{enumerate}
\def\labelenumi{(\alph{enumi})}
\tightlist
\item
  证明：对于每一个交换A-代数B，该合成法则将每一个有序对
\end{enumerate}

\[
(u, v) \in \operatorname{Hom} _ {\mathbf {A} - \mathrm{alg.}} (\mathbf {E}, \mathbf {B}) \times \operatorname{Hom} _ {\mathbf {A} - \mathrm{alg.}} (\mathbf {E}, \mathbf {B})
\]

对应到 A-代数同态

\[
\mathrm{E} \xrightarrow {c} \mathrm{E} \otimes \mathrm{E} \xrightarrow {u \otimes v} \mathrm{B} \otimes \mathrm{B} \xrightarrow {m _ {\mathrm{B}}} \mathrm{B}
\]

（其中 \(m_{\mathbf{B}}\) 是 B 中的乘积）在 \(\operatorname{Hom}_{\mathbf{A}\text{-alg.}}(\mathbf{E}, \mathbf{B})\) 上定义了一个幺半群结构。（b）为使对每个交换 A-代数 B， \(\operatorname{Hom}_{\mathbf{A}\text{-alg.}}(\mathbf{E}, \mathbf{B})\) 都是一个群，必要且充分的是存在一个 A-代数同态 \(i: \mathbf{E} \to \mathbf{E}\) ，使得 \(i(e) = e\) ，且复合映射 \(m \circ (1_{\mathbf{E}} \otimes i) \circ c\) 与 \(m \circ (i \otimes 1_{\mathbf{E}}) \circ c\) 都等于映射 \(x \mapsto \gamma(x)e\) 。这时称 \(i\) 为 E 中的一个逆；它是从 E 到相反双代数 \(E^0\) 上的一个同构。

\begin{enumerate}
\def\labelenumi{\arabic{enumi}.}
\setcounter{enumi}{4}
\item
  设 \(E_{1}\) , \(E_{2}\) 是 A 上的两个双代数。证明在张量积 \(E_{1} \otimes_{A} E_{2}\) 上， \(E_{1}\) 和 \(E_{2}\) 上的代数与余代数结构的张量积在 \(E_{1} \otimes_{A} E_{2}\) 上定义了一个双代数结构。
\item
  设 M 是类型为 N 的分次 A-模。在泛反交换代数 G(M) 上定义一个斜分次双代数结构（§ 7，习题
\end{enumerate}

13）；由此导出一个典范的分次代数同态 \(\mathbf{G}(\mathbf{M}^{*}) \to \mathbf{G}(\mathbf{M})^{*\mathbf{gr}}\) ，以及 \(\mathbf{G}(\mathbf{M})\) 的元素与 \(\mathbf{G}(\mathbf{M}^{*})\) 的元素之间的内积概念。

\begin{enumerate}
\def\labelenumi{\arabic{enumi}.}
\setcounter{enumi}{6}
\tightlist
\item
  在命题 11（第 9 号）的假设和记号下，证明公式
\end{enumerate}

\[
(\det g) \phi^ {- 1} = \bigwedge (g) \circ \phi^ {- 1} \circ \bigwedge (^ {t} g).
\]

由该公式与 (79)（第 11 号）可导出用余子式矩阵表示方阵的逆的表达式的一个新证明（§ 8，第 6 号，公式 (26)）。

\begin{enumerate}
\def\labelenumi{\arabic{enumi}.}
\setcounter{enumi}{7}
\tightlist
\item
  设 E 是维数为 n 的自由 A-模。~对从 \(\bigwedge(\mathbf{M})\) 到一个 A-模 G 中的每个 A-线性映射 v，我们可以写成 \(v = \sum_{p=0}^{n} v_{p}\) ，其中 \(v_{p}\) 是 v 在 \(\bigwedge^{p}(\mathbf{M})\) 上的限制；然后我们设 \(\eta v = \sum_{p=0}^{n} (-1)^{p} v_{p}\) ，使得 \(\eta^{2} v = v\) 。设 u 是从 M 到其对偶 \(M^{*}\) 上的一个同构，并设 \(\Delta\) 是 u 关于 M 的一组基 \((e_{i})\) 及 \(M^{*}\) 的对偶基的矩阵的行列式（该行列式不依赖于 M 的基的选取）。若 \(\phi\) 是从 \(\bigwedge(\mathbf{M})\) 到 \(\bigwedge(\mathbf{M}^{*})\) 上的、由 \(e = e_{1} \wedge e_{2} \wedge \cdots \wedge e_{n}\) 出发定义的同构，证明公式
\end{enumerate}

\[
\eta^ {n + 1} \Delta \phi = \bigwedge (t u) \circ \phi^ {- 1} \circ \bigwedge (u).
\]

\begin{enumerate}
\def\labelenumi{\arabic{enumi}.}
\setcounter{enumi}{8}
\tightlist
\item
  方阵 \(X\) （\(n\) 阶，在 \(A\) 上）的伴随矩阵是矩阵 \(\tilde{X} = (\det(X^{it}))\) （由 \(X\) 的 \(n - 1\) 阶子式组成）。证明若 \(X\) 可逆，则 \(\det(\tilde{X}) = (\det X)^{n-1}\) ，且每个子式 \(\det(\tilde{X}_{H,K})\) （\(\tilde{X}\) 的 \(p\) 阶子式）由公式给出
\end{enumerate}

\[
\det (\tilde {X} _ {\mathrm{H}, \mathrm{K}}) = (\det X) ^ {p - 1} (X _ {\mathrm{H} ^ {\prime}, \mathrm{K} ^ {\prime}})\tag{1}
\]

其中 H' 和 K' 分别是 H 和 K 在 \([1, n]\) 中的补集（“雅可比恒等式”；利用第 11 号的关系式 (80)）。

\begin{enumerate}
\def\labelenumi{\arabic{enumi}.}
\setcounter{enumi}{9}
\tightlist
\item
  从每个恒等式 \(\Phi = 0\) （任意可逆方阵 \(X\) （\(n\) 阶，在 A 上）的子式之间的）可以导出另一个恒等式 \(\tilde{\Phi} = 0\) ，称为 \(\Phi = 0\) 的余恒等式，其做法是将恒等式 \(\Phi = 0\) 应用于伴随矩阵 \(\tilde{X}\) （\(X\) 的）的子式，然后利用习题 9 中的恒等式 (1) 把 \(\tilde{X}\) 的子式替换为 \(X\) 的子式的函数。用这种方法证明以下恒等式：
\end{enumerate}

\[
\det (X ^ {i h}) \det (X ^ {j k}) - \det (X ^ {i k}) \det (X ^ {j h}) = (\det X) \det (X ^ {i j, h k}) \quad \text { for } \quad i <   j
\]

其中 \(X^{ij, hk}\) 表示 X 的 n-2 阶子式，它是通过在 X 中删去指标为 i、j 的行以及指标为 h、k 的列而得到的。

¶ 11. 设 \(\Phi = 0\) 是交换环 A 上一个可逆方阵 \(X\) （\(n\) 阶）的子式之间的恒等式， \(\tilde{\Phi} = 0\) 为其余恒等式（习题 10），而 \(k\) 为一个整数 \(>0\) 。设 \(Y\) 是 \(n + k\) 阶可逆方阵，并设 \(\tilde{Y}_0\) 是伴随矩阵 \(\tilde{Y}\) 通过在 \(\tilde{Y}\) 中删去指标 \(\leqslant k\) 的行和指标 \(\leqslant k\) 的列所成的子矩阵。若假设 \(\tilde{Y}_0\) 可逆，并将恒等式 \(\tilde{\Phi} = 0\) 应用于 \(\tilde{Y}_0\) 的子式，然后利用习题 9 的恒等式 (1)，把该恒等式中出现的每个子式（视为 \(\tilde{Y}\) 的子式）替换为用 \(Y\) 的子式表示的表达式，就得到一个恒等式 \(\Phi_k = 0\) （在矩阵 \(Y\) 的子式之间，当 \(Y\) 和 \(\tilde{Y}_0\) 可逆时成立），称为 \(k\) 阶的扩张，即恒等式 \(\Phi = 0\) 的扩张。

特别地，设 \(A = (\alpha_{ij})\) 是一个 \(n + k\) 阶可逆方阵，B 是 A 的 k 阶子矩阵（在 A 中删去指标 \textgreater k 的行和列而得），而 \(\Delta_{ij}\) 是 \(k + 1\) 阶矩阵的行列式（在 A 中删去指标 \textgreater k 但保留指标 \(k + i\) 的行、以及指标 \textgreater k 但保留指标 \(k + j\) 的列而得）；若 C 表示 n 阶矩阵 \((\Delta_{ij})\) ，证明恒等式

\[
\det C = (\det A) (\det B) ^ {n - 1}
\]

当 \(A\) 和 \(B\) 可逆时成立（证明它是 \(k\) 阶的扩张，施于 \(n\) 阶行列式的完全展开）。

叙述通过扩张拉普拉斯展开（§ 8，第 6 号）以及 § 8 习题 2 的恒等式所得到的恒等式。

¶ 12. 设 \(X = (\xi_{ij})\) 是一个 \(n\) 阶方阵（在交换域 \(K\) 上）；令 \(H\) 是 \([1, n]\) 的具有 \(p\) 个元素的一个子集，而 \(H'\) 为 \(H\) 在 \([1, n]\) 中的补集；假设对每对有序指标 \(h \in H\) , \(k \in H'\) ，

\[
\sum_ {i} \xi_ {i h} \xi_ {i k} = 0.
\]

证明：对 \([1, n]\) 中具有 p 个元素的所有子集 L, M ，有

\[
\rho_ {\mathbf {M}, \mathbf {M} ^ {\prime}} \det (X _ {\mathrm{L}, \mathrm{H}}) \det (X _ {\mathbf {M} ^ {\prime}, \mathbf {H} ^ {\prime}}) - \rho_ {\mathrm{L}, \mathrm{L} ^ {\prime}} \det (X _ {\mathbf {M}, \mathbf {H}}) \det (X _ {\mathrm{L}, \mathrm{H} ^ {\prime}}) = 0.
\]

（把列 \(x_h\) （指标 \(h \in \mathbf{H}\) ，在 \(X\) 中）看作 \(\mathbf{E} = \mathbf{K}^n\) 中的向量，把列 \(x_k^*\) （指标 \(k \in \mathbf{H}'\) ）看作 \(\mathbf{E}^*\) 中的向量，并证明 \((n - p)\) -形式 \(x_{\mathbf{H}'}^*\) 与 \(\phi_p(x_{\mathbf{H}})\) 成比例。）

\begin{enumerate}
\def\labelenumi{\arabic{enumi}.}
\setcounter{enumi}{12}
\tightlist
\item
  设 \(\Gamma\) 和 \(\Delta\) 是交换环上两个 n 阶可逆矩阵的行列式，且 \(\Gamma_{ij}\) 是通过将 \(\Gamma\) 中第 i 列替换为 \(\Delta\) 的第 j 列所得到的行列式。证明
\end{enumerate}

\[
\det (\Gamma_ {i j}) = \Gamma^ {n - 1} \Delta .
\]

（把 \(\Gamma_{ij}\) 按第 \(i\) 列展开，并利用习题 9。）

\begin{enumerate}
\def\labelenumi{\arabic{enumi}.}
\setcounter{enumi}{13}
\item
  设 M 为维数 n \textgreater{} 1 的自由 A-模。证明：每个同构 \(\psi\) （从 \(\bigwedge^{p}(M)\) 到 \(\bigwedge^{n-p}(M^{*})\) 上），使得 \(\bigwedge^{n-p}(\check{u}) \circ \psi = \psi \circ \bigwedge^{p}(u)\) 对 M 的每个行列式等于 1 的自同构 u 成立，都是第 11 号命题 12 中定义的同构 \(\phi_{p}\) 之一（按 § 7 习题 4 的方法进行）。
\item
  设 E 为维数为 n 的自由 A-模。~设 z 为 E 上的一个 p-向量， \(z^{*}\) 为 E 上的一个 \((p + q)\) -形式；z 可按典范方式（§ 7 习题 8）等同于 \(az_{1}\) （一个 p 阶反变张量 \(z_{1}\) 的斜对称化），而 \(z^{*}\) 等同于一个 \(p + q\) 阶协变张量的斜对称化。如果在混合张量 \(z_{1}\) . \(z^{*}\) 中，第 k 个反变指标与第 \((p + k)\) 个协变指标对 \(1 \leqslant k \leqslant p\) 进行缩并，证明如此得到的 q 阶协变张量是斜对称化的，并且除一个符号外典范地等同于 q-形式 \(z^{*} \perp z\) 。
\item
  设 E 是维数为 n 的自由 A-模。~ \(x \in \bigwedge(\mathrm{E})\) 与 \(y \in \bigwedge(\mathrm{E})\) 关于一个基 \(\{e\}\) （属于 \(\bigwedge^{n}(\mathrm{E})\) ）的回归积（记作 \(x \vee y\) ）是元素 \(\phi(\phi(x) \wedge \phi(y))\) ，同构 \(\phi\) 是相对于 e 的。该乘积仅在相差一个可逆因子的意义下有定义，该因子依赖于所选的基 \(\{e\}\) 。证明：若 x 是一个 p-向量，y 是一个 q-向量，则 \(x \vee y = 0\) 当 \(p + q < n\) ，且若 \(p + q \geqslant n\) ， \(x \vee y\) 是一个 \((p + q - n)\) -向量，使得
\end{enumerate}

\[
y \vee x = (- 1) ^ {(n - p) (n - q)} x \vee y.
\]

回归积是结合的、对加法是分配的，并在 \(\bigwedge(\mathrm{E})\) 上定义了一个含单位元的、同构于外代数结构的代数结构。把 \(x \vee y\) 的分量表示为 \(x\) 和 \(y\) 的分量（对 \(\mathbf{E}\) 的一个给定基）的函数。

\begin{enumerate}
\def\labelenumi{\arabic{enumi}.}
\setcounter{enumi}{16}
\tightlist
\item
  设 K 是一个交换域，E 是 K 上的一个向量空间。
\end{enumerate}

\begin{enumerate}
\def\labelenumi{(\alph{enumi})}
\item
  设 \(z\) 是一个 \(p\) -向量 \(\neq 0\) （在 \(E\) 上），并设 \(V_{z}\) 为 \(E\) 的由向量 \(x\) （满足 \(z \wedge x = 0\) ）组成的向量子空间。则 \(\dim(V_{z}) \leqslant p\) ；若 \((x_{i})_{1 \leqslant i \leqslant q}\) 是 \(V_{z}\) 中向量的一个自由系，则存在一个 \((p - q)\) -向量 \(v\) ，使得 \(z = v \wedge x_{1} \wedge \dots \wedge x_{q}\) 。若 \(E\) 是有限维的，则 \(V_{z} = \phi(M_{z}^{0})\) ，其中 \(M_{z}^{0}\) 是 \(E^{*}\) 中的子空间 \(M_{z}\) （与 \(z\) 相关联）的正交补。要使 \(z\) 为纯 \(p\) -向量，必要且充分的是 \(\dim(V_{z}) = p\) ，此时 \(V_{z} = M_{z}\) 。
\item
  设 \(v\) 为纯 \(p\) -向量， \(w\) 为纯 \(q\) -向量，而 \(V\) 和 \(W\) 为 \(E\) 的分别与 \(v\) 和 \(w\) 相关联的子空间。要使 \(V \subset W\) ，必要且充分的是 \(q \geqslant p\) 且存在一个 \((q - p)\) -向量 \(u\) ，使得 \(w = u \wedge v\) 。要使 \(V \cap W = \{0\}\) ，必要且充分的是 \(v \wedge w \neq 0\) ；此时子空间 \(V + W\) 与纯 \((p + q)\) -向量 \(v \wedge w\) 相关联。
\item
  设 \(u\) 和 \(v\) 为两个纯 \(p\) -向量，且 \(U\) 和 \(V\) 为分别与 \(u\) 和 \(v\) 相关联的子空间。要使 \(u + v\) 为纯的，必要且充分的是 \(\dim(U \cap V) \geqslant p - 1\) 。
\end{enumerate}

\begin{enumerate}
\def\labelenumi{\arabic{enumi}.}
\setcounter{enumi}{17}
\item
  设 \(z = \sum_{H} \alpha_{H} e_{H}\) 为 n 维向量空间 E 的一个非零纯 p-向量，用它关于 E 的任意基 \((e_{i})_{1 \leqslant i \leqslant n}\) 的分量表示。设 G 为 \([1, n]\) 的具有 p 个元素的子集，使得 \(\alpha_{G} \neq 0\) ；令 \((i_{h})_{1 \leqslant h \leqslant p}\) 为 G 中指标按递增顺序排列的序列，而 \((j_{k})_{1 \leqslant k \leqslant n - p}\) 为 G 的补集 \(G'\) 中指标按递增顺序排列的序列。对每对有序指标 \((h, k)\) （满足 \(1 \leqslant h \leqslant p\) , \(1 \leqslant k \leqslant n - p\) ），令 \(\beta_{hk}\) 为 z 的分量 \(\alpha_{H}\) （对应于 \([1, n]\) 的子集 H），该子集 H 由 G 中异于 \(i_{h}\) 的 p - 1 个指标以及指标 \(j_{k} \in G'\) 组成；设 X 为具有 p 行和 n - p 列的矩阵 \((\beta_{hk})\) 。给定 \([1, n]\) 的任意 p 元子集 L，使得 \(L \cap G\) 具有 \(q \geqslant 1\) 个元素，证明 \((\alpha_{G})^{q-1} \alpha_{L}\) 除符号外等于 X 的 q 阶子式，该子式由满足 \(i_{h} \in G \cap G\) 的指标 h 的行和满足 \(j_{k} \in L \cap G\) 的指标 k 的列组成。（把 z 写成 \(\alpha_{G} y_{1} \wedge y_{2} \wedge \cdots \wedge y_{p}\) 的形式，其中向量 \(y_{i}\) 使得：在以 \(y_{i}\) 为列的 n 行 p 列矩阵 Y 中，由指标属于 G 的行组成的子矩阵是单位矩阵。）
\item
  在有限维向量空间 E 上，设 \(z = x_{1} \wedge x_{2} \wedge \cdots \wedge x_{p}\) 为一个纯p-向量 \(\neq 0\) ；令 \(v^{*}\) 为一个q-形式（ \(q \leqslant p\) ) 且 \(u^{*}\) 的任意一个元素 \(\bigwedge(\mathrm{E}^{*})\) 。证明
\end{enumerate}

\[
(u ^ {*} \wedge v ^ {*}) \lrcorner z = (- 1) ^ {q (q - 1) / 2} \sum_ {\mathrm{H}} \rho_ {\mathrm{H}, \mathrm{K}} \langle v ^ {*}, x _ {\mathrm{H}} \rangle (u ^ {*} \lrcorner x _ {\mathrm{K}})
\]

其中 H 遍历所有子集的集合 \([1, p]\) 具有q个元素，K是H在 \([1, p]\) 和 \(x_{H}\) （相应地 \(x_{K}\) )表示 \(x_{i}\) 的集合，使得 \(i \in H\) （相应地 \(i \in K\) ).

\begin{enumerate}
\def\labelenumi{\arabic{enumi}.}
\setcounter{enumi}{19}
\item
  设 E 为 n 维向量空间，z 为 E 上的纯 p-向量，V 为 E 的与 z 相关联的向量子空间， \(u^{*}\) 为 \(E^{*}\) 上的一个纯 q-向量， \(W'\) 为 \(E^{*}\) 的与 \(u^{*}\) 相关联的子空间，而 W 为 E 的、维数为 n - q 且正交于 \(W'\) 的子空间。若 q \textless{} p，则 \(u^{*} \lrcorner z\) 是 E 上的一个纯 (p - q)-向量，它当 \(\dim(V \cap W) > p - q\) 时为零；若 \(\dim(V \cap W) = p - q\) ，则 \(V \cap W\) 与 \(u^{*} \lrcorner z\) 相关联。
\item
  \begin{enumerate}
  \def\labelenumii{(\alph{enumii})}
  \tightlist
  \item
    直接证明（不使用同构 \(\phi_p\) ）：每个 \((n - 1)\) -向量（在 \(n\) 维向量空间上的）都是纯的。
  \end{enumerate}
\end{enumerate}

\begin{enumerate}
\def\labelenumi{(\alph{enumi})}
\setcounter{enumi}{1}
\tightlist
\item
  设 A 是域 K 上的一个交换代数，其基由单位元 1 和三个元素 \(a_{1}\) , \(a_{2}\) , \(a_{3}\) 组成，它们彼此之间的乘积为零。设 E 为 A-模 \(\mathbf{A}^3\) ，而 \((e_i)_{1\leqslant i\leqslant 3}\) 为其典范基；证明双向量
\end{enumerate}

\[
a _ {1} (e _ {2} \land e _ {3}) + a _ {2} (e _ {3} \land e _ {1}) + a _ {3} (e _ {1} \land e _ {2})
\]

不是纯的。

\begin{enumerate}
\def\labelenumi{\arabic{enumi}.}
\setcounter{enumi}{21}
\tightlist
\item
  设 E 为一个有序集，其中每个区间 \([a, b]\) 都是有限的。设 S 为 \(E \times E\) 的由有序对 \((x, y)\) （满足 \(x \leqslant y\) ）组成的子集。证明在 A-模 \(\mathbf{A}^{(\mathrm{s})}\) （由 S 中元素的形式线性组合构成）上，通过取余积定义了一个余结合的余代数结构
\end{enumerate}

\[
c ((x, y)) = \sum_ {x \leqslant z \leqslant y} (x, z) \otimes (z, y).
\]

对于这个余代数，线性形式 \(\gamma\) 使得

\[
\gamma ((x, y)) = 0 \quad \text { if } \quad x \neq y \text { in   E }, \quad \gamma ((x, x)) = 1
\]

是一个余单位。

\begin{enumerate}
\def\labelenumi{\arabic{enumi}.}
\setcounter{enumi}{22}
\tightlist
\item
  设 C 是环 A 上的一个余代数。C 的一个子 A-模 S 称为 C 的子余代数，如果 \(c(S) \subset \operatorname{Im}(S \otimes S)\) 。
\end{enumerate}

\begin{enumerate}
\def\labelenumi{(\alph{enumi})}
\item
  若 C 是余结合的（相应地，余交换的），则其每个子余代数也是。若 \(\gamma\) 是 C 的一个余单位，则 \(\gamma\) 在 C 的任意子余代数上的限制是一个余单位。
\item
  余代数 \(\mathbf{A}^{(\mathbf{X})}\) （第 1 号，例 4）的子余代数有哪些？
\item
  若 \((\mathbf{S}_{\alpha})\) 是 \(\mathbf{C}\) 的一个子余代数族，则 \(\sum_{\alpha} \mathbf{S}_{\alpha}\) 是 \(\mathbf{C}\) 的一个子余代数（包含所有 \(\mathbf{S}_{\alpha}\) 的最小者）。
\end{enumerate}

\begin{enumerate}
\def\labelenumi{\arabic{enumi}.}
\setcounter{enumi}{23}
\tightlist
\item
  设 C 是环 A 上的一个余代数。C 的一个子 A-模 S 称为右（相应地，左）余理想，如果 \(c(S) \subset \operatorname{Im}(S \otimes C)\) （相应地， \(c(S) \subset \operatorname{Im}(C \otimes S)\) ）。
\end{enumerate}

\begin{enumerate}
\def\labelenumi{(\alph{enumi})}
\item
  子余代数（习题 23）既是右余理想又是左余理想。反之，若 \(\mathbf{A}\) 是一个域，则结论成立。
\item
  每个右（相应地，左）余理想的并都是右（相应地，左）余理想。
\item
  如果 S 是 C 的右（相应地，左）余理想，则正交补 \(S^{0}\) 是对偶代数 \(C^{*}\) 的一个右（相应地，左）理想。
\item
  若 A 是一个域，且 \(\mathfrak{I}'\) 是对偶代数 \(\mathbf{C}^*\) 的一个右（相应地，左）理想，则正交 \(\mathfrak{I}''^0\) （在 \(\mathbf{C}\) 中）是 \(\mathbf{C}\) 的一个右（相应地，左）余理想。（用反证法论证：假设存在 \(x \in \mathfrak{I}''^0\) ，使得 \(c(x) \notin \mathfrak{I}''^0 \otimes \mathbf{C}\) ；写成 \(c(x) = \sum_{i} y_i \otimes z_i\) ，其中 \(z_i\) 线性无关，并证明这将导致存在 \(y' \in \mathfrak{I}'\) 和 \(z' \in \mathbf{C}^*\) ，使得 \(\langle y' z', x \rangle \neq 0\) 。）
\end{enumerate}

由此推出：对于 C 的一个向量子空间 S，要使 S 成为右（相应地，左）余理想，当且仅当 \(S^{0}\) 是 \(C^{*}\) 的一个右（相应地，左）理想。

\begin{enumerate}
\def\labelenumi{(\alph{enumi})}
\setcounter{enumi}{4}
\tightlist
\item
  由 (d) 推出：若 A 是一个域，则 C 的任意一族右余理想（相应地，左余理想，相应地，子余代数）的交仍是右余理想（相应地，左余理想，相应地，子余代数）。要使 S 成为 C 的子余代数，必要且充分的是其正交补 \(S^{0}\) 是代数 \(C^{*}\) 的一个双边理想。
\end{enumerate}

\begin{enumerate}
\def\labelenumi{\arabic{enumi}.}
\setcounter{enumi}{24}
\tightlist
\item
  设 C 是环 A 上的一个余代数。C 的一个子 A-模 S 称为余理想，如果 \(c(\mathbf{S}) \subset \operatorname{Im}(\mathbf{S} \otimes \mathbf{C}) + \operatorname{Im}(\mathbf{C} \otimes \mathbf{S})\) 。若 C 是有余单位的，其余单位为 \(\gamma\) ，则余理想 S 称为余零的，如果 \(\gamma(x) = 0\) （对 \(x \in S\) ）。
\end{enumerate}

\begin{enumerate}
\def\labelenumi{(\alph{enumi})}
\item
  右或左余理想都是余理想。余理想的任意和仍是余理想。
\item
  若 S 是 C 的余理想，则其正交补 \(S^{0}\) 是对偶代数 \(C^{*}\) 的一个子代数。若 C 是有余单位的且 S 是余零的，则 \(S^{0}\) 是含单位元的。
\item
  设 A 是一个域，C 是一个有余单位的余代数。证明若 \(E'\) 是 \(C^{*}\) 的一个含单位元的子代数（与 \(C^{*}\) 有相同的单位元），则 C 中的正交 \(E'^{0}\) 是一个余零余理想。
\item
  如果 S 是 C 的一个余理想，证明在 C/S 上存在且仅存在一个余代数结构，使得典范映射 \(p: C \rightarrow C/S\) 是一个余代数态射。如果 C 是有余单位的且 S 是余零的，则 C/S 是有余单位的。
\item
  对于每个余代数态射 \(f\colon\mathbf{C}\to\mathbf{C}^{\prime},\mathbf{S}=\operatorname{Ker}(f)\) 是 \(\mathbf{C}\) 的余理想， \(\mathbf{S}^{\prime}=\operatorname{Im}(f)\) 是 \(\mathbf{C}^{\prime}\) 的一个子余代数，并且在典范分解
\end{enumerate}

\[
\mathrm{C} \xrightarrow {p} \mathrm{C} / \mathrm{S} \xrightarrow {g} \mathrm{S} ^ {\prime} \xrightarrow {j} \mathrm{C} ^ {\prime}
\]

中 \(f, g\) 是一个余代数同构。

\begin{enumerate}
\def\labelenumi{(\arabic{enumi})}
\setcounter{enumi}{25}
\tightlist
\item
  设 C 是交换域 K 上的一个余结合、有余单位的余代数。
\end{enumerate}

\begin{enumerate}
\def\labelenumi{(\alph{enumi})}
\item
  对于 C 的每一个向量子空间 E（它在合成律 \((u, x) \mapsto u \perp x\) 下是一个左 C*-模），这个 C*-模的左零化子是一个双边理想 \(S'\) （在代数 \(C^*\) 中）。证明若 E 在 K 上是有限维的，则 \(C^*/S'\) 是有限维的。
\item
  从 (a) 推出，对每个元素 \(x \in \mathbb{C}\) ， \(\mathbb{C}\) 的由 \(x\) 生成的子余代数是有限维的。（注意， \(\mathbb{C}^*\) -模 \(\mathbb{E}\) （由 \(x\) 生成的）是有限维的，并且若 \(\mathfrak{S}'\) 是其左零化子，则 \(x \in \mathfrak{S}'^0\) ，即 \(\mathfrak{S}'\) 在 \(\mathbb{C}\) 中的正交。）
\end{enumerate}

\begin{enumerate}
\def\labelenumi{(\arabic{enumi})}
\setcounter{enumi}{26}
\tightlist
\item
  设 B 是交换域 K 上的一个结合含幺代数，并设 \(m: B \otimes_{K} B \to B\) 是定义 B 上乘法的 K-线性映射。张量积 \(B^{*} \otimes_{K} B^{*}\) （其中 \(B^{*}\) 是向量空间 B 的对偶向量空间）典范地等同于 \((\mathbf{B} \otimes_{\mathbf{K}} \mathbf{B})^{*}\) 的一个向量子空间（II，§7，第 7 号，命题 16）。
\end{enumerate}

\begin{enumerate}
\def\labelenumi{(\alph{enumi})}
\item
  证明对于元素 \(w \in \mathbf{B}^*\) ，以下条件等价：（1） \({}^t m(w) \in \mathbf{B}^* \otimes \mathbf{B}^*\) ；(2) 元素 \(x \lrcorner w\) （其中 \(x\) 遍历 \(\mathbf{B}\) ）组成的集合是 \(\mathbf{B}^*\) 的一个有限维子空间；(3) 元素 \(w \lfloor x\) （其中 \(x\) 遍历 \(\mathbf{B}\) ）组成的集合是 \(\mathbf{B}^*\) 的一个有限维子空间；(4) 元素 \(x \lrcorner w \lfloor y\) （其中 \(x\) 和 \(y\) 遍历 \(\mathbf{B}\) ）组成的集合包含在 \(\mathbf{B}^*\) 的一个有限维子空间中。
\item
  设 \(\mathbf{B}' \subset \mathbf{B}^*\) 为 \(w \in \mathbf{B}^*\) 中满足
\end{enumerate}

(a) 中的等价条件者的集合。证明 \({}^t m(\mathbf{B}') \subset \mathbf{B}' \otimes_{\mathbb{K}} \mathbf{B}'\) 。（对于 \(w \in \mathbf{B}'\) ，记 \({}^t m(w) = \sum_i u_i \otimes v_i\) ，其中例如 \(u_i\) 在 \(\mathbf{B}^*\) 中线性无关；然后证明 \(v_i \in \mathbf{B}'\) （对所有 \(i\) ），利用 \(\mathbf{B}\) 的结合性并通过反证法论证。）\(\mathbf{B}'\) 因此是包含在 \(\mathbf{B}^*\) 中的最大余代数，对其而言 \({}^t m\) 是余积。\(\mathbf{B}'\) 称为代数 \(\mathbf{B}\) 的对偶余代数。当 \(\mathbf{B}\) 是有限维的时， \(\mathbf{B}' = \mathbf{B}^*\) 。

\begin{enumerate}
\def\labelenumi{(\alph{enumi})}
\setcounter{enumi}{2}
\item
  证明：如果 \(\mathbf{B}\) 是 \(\mathbf{K}\) 上无限维的交换域，则 \(\mathbf{B}' = \{0\}\) 。（注意，元素 \(x \perp w\) （对固定的 \(w \in \mathbf{B}^*\) ，\(x\) 遍历 \(\mathbf{B}\) ）组成的集合的正交补是 \(\mathbf{B}\) 的一个理想。）
\item
  设 \(\mathbf{B}_1, \mathbf{B}_2\) 是两个 K-代数，且 \(f: \mathbf{B}_1 \to \mathbf{B}_2\) 是代数的一个 K-同态。证明 \(^t f(\mathbf{B}_2') \subset \mathbf{B}_1'\) ，并且 \(^t f\) 限制在 \(\mathbf{B}_2'\) 上是从 \(\mathbf{B}_2'\) 到 \(\mathbf{B}_1'\) 中的余代数同态。
\item
  若 \(\mathbf{C}\) 是 \(\mathbf{K}\) 上的一个余结合、有余单位的余代数，则典范单射 \(\mathbf{C} \to \mathbf{C}^{**}\) （向量空间 \(\mathbf{C}\) 到其二对偶中的）把 \(\mathbf{C}\) 映到 \((\mathbf{C}^*)'\) ——代数 \(\mathbf{C}^*\) （对偶于 \(\mathbf{C}\) ）的对偶余代数——的一个子空间上，并且是从 \(\mathbf{C}\) 到 \((\mathbf{C}^*)'\) 中的余代数同态。
\end{enumerate}

